\section{Results}

We report results in the order of the evaluation tiers of Section~3.2: first the
seven-model pre-run observations (image reliance, degradation, and failure
modes), then the two-primary-model NIH pilot and shift benchmark, then the
external report evaluation. Pre-run quantities are point estimates on 100 images
per model and are labelled as such wherever they appear.

\subsection{Structured outputs and parse stability}

The two primary models produced structurally valid JSON for all 1,000 NIH pilot studies after the prompt and parser were finalised (Table~\ref{tab:parse}). Parse stability, however, did not imply reliable image understanding, as the following subsections show.

\begin{table}[t]
\centering
\small
\begin{tabular}{lll}
\hline
Model & Samples & Parse failures \\
MedGemma-4B-IT & 1000 & 0 \\
Qwen2.5-VL-7B & 1000 & 0 \\
\hline
\end{tabular}
\caption{Structured output parse stability on the NIH pilot.}
\label{tab:parse}
\end{table}

\subsection{Seven-model image-reliance spectrum}

Across all seven models, replacing the image with a blank produces a graded
change in the models' positive probabilities, from a large change to none at all
(Table~\ref{tab:reliance-spectrum}). Three tiers emerge mechanically: two models
that read the image (visual-reliance signal $\geq 0.33$), two that are partially
dependent (0.13--0.16), and two that are effectively image-blind ($\leq 0.012$),
with CheXagent's probabilities byte-stable across every degenerate input. Every
model is more sensitive to the blank than to blur or shuffle, which is what a
model that reads pixels should do; the image-blind models show no such ordering.

\begin{table}[t]
\centering
\small
\begin{tabular}{lcccc}
\hline
Model & VRS blank & VRS blur & VRS shuffle & Status flip (blank) \\
\hline
InternVL3.5-8B & 0.363 & 0.248 & 0.262 & 0.393 \\
MedGemma-4B-IT & 0.334 & 0.210 & 0.187 & 0.498 \\
MedGemma-1.5-4B-IT (bf16) & 0.283 & 0.235 & 0.269 & 0.511 \\
Qwen3-VL-8B-Instruct & 0.156 & 0.163 & 0.241 & 0.189 \\
MedGemma-1.5-4B-IT (4-bit) & 0.128 & 0.120 & 0.118 & 0.228 \\
Qwen2.5-VL-7B-Instruct & 0.012 & 0.012 & 0.028 & 0.015 \\
CheXagent-8B & 0.000 & 0.000 & 0.000 & 0.000 \\
\hline
\end{tabular}
\caption{Visual-reliance signal (VRS) and blank-image status-flip rate on the
seven-model pre-run (n=100 images per model). Models are ordered by VRS under the
blank condition.}
\label{tab:reliance-spectrum}
\end{table}

The gradient is not a binary split. The two within-family pairs refine it:
MedGemma-1$\rightarrow$1.5 keeps the family image-dependent, while
Qwen2.5-VL$\rightarrow$Qwen3-VL moves a family from image-blind to partially
dependent, so generational change can alter image reliance within a family. A
precision interaction is visible on the one model with both a 4-bit and a bf16
row (MedGemma-1.5): the 4-bit VRS (0.128) is less than half the bf16 VRS
(0.283), an observed association on a single model rather than a controlled
ablation.

\subsection{Two primary models on the NIH pilot}

Macro metrics across the eight findings for the two primary models are shown in
Table~\ref{tab:main}.

\begin{table}[t]
\centering
\small
\begin{tabular}{lllllll}
\hline
Model & Split & Macro AUROC & Macro ECE & Macro Brier & Macro AURC & Mean coverage @5\% risk \\
MedGemma & internal\_train & 0.651 & 0.129 & 0.119 & 0.120 & 0.733 \\
MedGemma & source\_val & 0.627 & 0.147 & 0.119 & 0.137 & 0.713 \\
MedGemma & internal\_test & 0.622 & 0.220 & 0.188 & 0.228 & 0.357 \\
Qwen2.5-VL & internal\_train & 0.500 & 0.088 & 0.087 & 0.071 & 0.662 \\
Qwen2.5-VL & source\_val & 0.502 & 0.070 & 0.070 & 0.061 & 0.662 \\
Qwen2.5-VL & internal\_test & 0.509 & 0.153 & 0.154 & 0.117 & 0.398 \\
\hline
\end{tabular}
\caption{Macro performance and calibration on NIH data.}
\label{tab:main}
\end{table}

MedGemma showed moderate but non-trivial discrimination (macro AUROC 0.62--0.65)
and clear calibration degradation on \texttt{internal\_test} (macro ECE 0.220).
Qwen2.5-VL remained close to chance discrimination (macro AUROC $\approx$ 0.50),
despite acceptable-looking calibration and selective metrics.

On 300 NIH images, MedGemma changed many predictions under counterfactual image
conditions while Qwen2.5-VL changed almost none (Table~\ref{tab:image-reliance},
Figure~\ref{fig:image-reliance}).

\begin{table}[t]
\centering
\small
\begin{tabular}{lll}
\hline
Condition & MedGemma image sensitivity & Qwen2.5-VL image sensitivity \\
blank gray & 1.000 & 0.0033 \\
strong blur & 0.947 & 0.0033 \\
shuffled image & 0.877 & 0.0033 \\
\hline
\end{tabular}
\caption{Image reliance on 300 NIH studies.}
\label{tab:image-reliance}
\end{table}

Qwen2.5-VL predicted \texttt{no\_finding = positive} for 988/1000 studies and had
macro AUROC $\approx$ 0.50 across all eight findings. Its apparently low
selective risk and high selective accuracy therefore reflected a degenerate
prediction strategy rather than reliable image interpretation.

\textbf{Finding:} Structured output validity, low selective risk, and high selective
accuracy can coexist with image invariance.

\subsection{Shift effects are heterogeneous and non-monotonic in accuracy}

MedGemma macro AUROC and ECE under shift are shown in Table~\ref{tab:shift}.

\begin{table}[t]
\centering
\small
\begin{tabular}{llll}
\hline
Condition & Severity & Macro AUROC & Macro ECE \\
clean & none & 0.620 & 0.165 \\
blur & mild & 0.612 & 0.167 \\
blur & moderate & 0.603 & 0.164 \\
blur & severe & 0.615 & 0.142 \\
noise & mild & 0.675 & 0.155 \\
noise & moderate & 0.704 & 0.155 \\
noise & severe & 0.520 & 0.145 \\
JPEG & mild & 0.591 & 0.169 \\
JPEG & moderate & 0.710 & 0.153 \\
JPEG & severe & 0.621 & 0.148 \\
brightness & mild & 0.614 & 0.175 \\
brightness & moderate & 0.600 & 0.164 \\
brightness & severe & 0.597 & 0.166 \\
contrast & mild & 0.670 & 0.165 \\
contrast & moderate & 0.620 & 0.166 \\
contrast & severe & 0.617 & 0.166 \\
resolution & mild & 0.633 & 0.152 \\
resolution & moderate & 0.705 & 0.138 \\
resolution & severe & 0.581 & 0.109 \\
\hline
\end{tabular}
\caption{MedGemma macro AUROC and ECE under image-quality shift.}
\label{tab:shift}
\end{table}

Image sensitivity generally increased with severity (e.g., noise 0.53 $\rightarrow$ 0.70 $\rightarrow$ 0.79; resolution 0.37 $\rightarrow$ 0.60 $\rightarrow$ 0.71), but AUROC did not follow a monotonic degradation pattern. Several shifts even improved calibration. Paired macro bootstrap tests on MedGemma showed no shift produced a significant AUROC decrease; JPEG moderate ($\Delta$AUROC = +0.069, p = 0.05) and resolution moderate ($\Delta$AUROC = +0.068, p = 0.10) were nominally positive, and several conditions significantly improved ECE, while only brightness mild significantly increased ECE ($\Delta$ECE = +0.010, p = 0.04).

\textbf{Finding:} Distribution shift does not produce a simple monotonic degradation
of CXR VLM reliability as measured by macro accuracy or calibration. Shift response
is condition-specific.

\subsection{Degradation suppresses positive recall, not macro accuracy}

Macro accuracy hides the clinically meaningful effect of shift. Across the seven
parse-complete models in the pre-run, macro accuracy stays flat through the shift
cells ($\approx$0.59--0.72, within $\pm$0.03 between clean and severe). The
flatness is an accuracy paradox: the pathology labels are minority classes (236
positive instances per cell), so a model that defaults to ``negative everywhere''
scores a respectable accuracy. The signal sits in the positive recall of the
pathology labels, where the dose-response is clear and degradation-type-specific
(Table~\ref{tab:recall}).

\begin{table}[t]
\centering
\small
\begin{tabular}{lcccc}
\hline
Model & Recall clean & Resolution severe & Noise severe & Blur severe \\
\hline
MedGemma-1.5-4B-IT & 0.64 & 0.46 & 0.41 & 0.60 \\
MedGemma-4B-IT & 0.59 & 0.34 & 0.32 & 0.53 \\
InternVL3.5-8B & 0.46 & 0.26 & 0.43 & 0.43 \\
Qwen3-VL-8B-Instruct & 0.21 & 0.13 & 0.13 & 0.09 \\
Qwen2.5-VL-7B-Instruct & 0.02 & 0.00 & 0.00 & 0.00 \\
CheXagent-8B & 0.00 & 0.04 & 0.05 & 0.01 \\
\hline
\end{tabular}
\caption{Pathology positive recall under the three counterfactual and three
severity conditions, pre-run (n=100). Resolution and noise suppress recall across
all image-reading models; blur leaves it near the clean level.}
\label{tab:recall}
\end{table}

Three observations follow. First, resolution and noise are the two degradations
that cut recall across every image-reading model (MedGemma-1.5, MedGemma-4B,
InternVL3.5, Qwen3-VL), while blur, brightness, and contrast leave recall near the
clean level. Resolution and noise are the clinically relevant pair -- JPEG-style
downsampling and low-dose photon noise. Second, the dose-response is monotone in
severity within those conditions (e.g.\ MedGemma-4B recall 0.59 $\rightarrow$ 0.50
$\rightarrow$ 0.34 down the resolution ladder). Third, Qwen2.5-VL and CheXagent have
no recall to lose: their failure is already complete on clean images, so degradation
cannot make it worse.

This is a pre-run observation, not a fitted production claim; the full run
re-censuses these curves and is the arbiter of whether the degradation-type
ranking (resolution $\approx$ noise $>$ JPEG $>$ blur/brightness/contrast)
holds at scale. We report it because it is parameterised and falsifiable, and
because macro accuracy -- the quantity a majority-reporting model excels at --
would not have exposed it.

\textbf{Finding:} On an imbalanced benchmark, shift degradation concentrates in
positive recall rather than macro accuracy, and the effect is
degradation-type-specific.

\subsection{An image-blind, default-normal model}

Three independent quantities place Qwen2.5-VL at an unsafe corner: its VRS is
0.012 (Table~\ref{tab:reliance-spectrum}), its pathology recall is $\approx$0.02
(Table~\ref{tab:recall}), and its \texttt{no\_finding} report rate rises from 0.930
on clean images to 0.990 at severe degradation (Table~\ref{tab:nofind}). The model
does not read the image and it reports ``no finding'' almost always, \emph{more}
often as the image degrades -- the hardest cases are the ones it most confidently
declares normal.

\begin{table}[t]
\centering
\small
\begin{tabular}{lcccc}
\hline
Model & Clean & Mild & Moderate & Severe \\
\hline
Qwen2.5-VL-7B-Instruct & 0.930 & 0.940 & 0.973 & 0.990 \\
Qwen3-VL-8B-Instruct & 0.649 & 0.638 & 0.609 & 0.642 \\
InternVL3.5-8B & 0.200 & 0.217 & 0.215 & 0.237 \\
MedGemma-4B-IT & 0.010 & 0.008 & 0.008 & 0.015 \\
MedGemma-1.5-4B-IT & 0.000 & 0.000 & 0.003 & 0.025 \\
CheXagent-8B & 0.000 & 0.000 & 0.000 & 0.000 \\
\hline
\end{tabular}
\caption{Probability of reporting \texttt{no\_finding} across shift severity,
pre-run (n=100). Qwen2.5-VL is the only model whose default-normal rate rises with
degradation.}
\label{tab:nofind}
\end{table}

The claim carried forward is the compound -- low visual reliance \emph{and} a
high default-normal rate in the same model -- because either alone is ambiguous.
A model with low VRS but balanced predictions, or one that defaults normal while
still reading the image, would not occupy the same unsafe corner.

\subsection{Failure modes without a parse error}

The seven models separate into three failure modes, none of which requires a
parse failure. Qwen2.5-VL is default-normal (Section~4.6). LLaVA-Med answers in
prose and never honours the structured contract, so its rows carry no parseable
JSON despite a well-formed pipeline. CheXagent is the refusing mode: on the shift
stage it emits \texttt{positive} exactly 8 times out of 15,200 label slots,
\texttt{not\_mentioned} on 77\% of slots, and never reports
\texttt{no\_finding = positive}. Its bare-label reply is effectively ``there is
something wrong, but I will not say what'' -- mechanically the opposite of
Qwen2.5-VL, which fills every slot and says ``healthy.''

\begin{table}[t]
\centering
\small
\begin{tabular}{lccccc}
\hline
Model & Positive & Negative & Uncertain & Not mentioned & No-finding rate \\
\hline
CheXagent-8B & 8 & 2{,}902 & 419 & 11{,}711 (77\%) & 0.000 \\
Qwen2.5-VL-7B-Instruct & 1{,}912 & 13{,}288 & 0 & 0 & 0.930--0.990 \\
\hline
\end{tabular}
\caption{Label-slot usage on the shift stage, pre-run (n=100). CheXagent refuses
most slots; Qwen2.5-VL fills them with a default normal.}
\label{tab:slot}
\end{table}

The three failure shapes -- default-normal, refusal-to-answer, and
prose-without-the-contract -- are a taxonomy the pipeline distinguishes without
resorting to parse-error penalties, and each invites a different clinical guard.

\subsection{Selective prediction}

On pooled \texttt{internal\_test} label-instances, MedGemma achieved a macro AURC
of 0.201 and coverage at 5\% risk of 0.002 by the fitted temperature scaler,
while Qwen2.5-VL achieved a lower AURC (0.095) but higher coverage at 10\% risk
(0.897) and higher selective accuracy at 80\% coverage (0.901). Qwen2.5-VL's
selective metrics appear strong but are driven by image-invariant majority
prediction rather than visual understanding, so they are not valid as clinical
evidence. Selective risk must therefore be conditioned on image reliance.

\subsection{CheXpert-plus-RRG external report evaluation}

On 264 RRG images, MedGemma produced an overall hallucination rate of 0.135, a
negation-failure rate of 0.006, and an omission rate of 0.122
(Table~\ref{tab:rrg}, Figure~\ref{fig:rrg-errors}).

\begin{table}[t]
\centering
\small
\begin{tabular}{lll}
\hline
Finding & Hallucination & Omission \\
atelectasis & 0.443 & 0.057 \\
cardiomegaly & 0.449 & 0.003 \\
consolidation & 0.019 & 0.127 \\
edema & 0.051 & 0.166 \\
no\_finding & 0.083 & 0.252 \\
pleural\_effusion & 0.019 & 0.213 \\
pneumonia & 0.019 & 0.054 \\
pneumothorax & 0.000 & 0.105 \\
\hline
\end{tabular}
\caption{Finding-level report errors on CheXpert-plus-RRG.}
\label{tab:rrg}
\end{table}

The high atelectasis/cardiomegaly hallucination rates may partly reflect model
over-calling and limitations of the rule-based evidence extractor. On the first
50 RRG images, the original-vs-blank status-flip rate was 0.4525, indicating that
MedGemma still depends on the image in this external set.

\subsection{Summary of empirical findings}

\begin{enumerate}
\item Visual reliance grades continuously across seven models, from image-dependent to image-blind.
\item An image-blind model can emit valid structured outputs and apparently favourable selective metrics.
\item Degradation concentrates in positive recall rather than macro accuracy, and is degradation-type-specific (resolution and noise dominate).
\item An image-blind model can compound with a default-normal strategy, so the hardest cases are declared normal.
\item Failure modes (default-normal, refusal, prose-without-contract) exist without any parse error.
\item Image-quality shift does not monotonically degrade macro accuracy or calibration.
\item External report evaluation exposes both hallucination and omission behaviour.
\end{enumerate}